\documentclass[11pt]{article}
\usepackage[a4paper,top=0.6cm,bottom=3.0cm,left=1.9cm,right=1.9cm,footskip=2.0cm]{geometry}
\usepackage{graphicx}
\usepackage[table]{xcolor}
\usepackage{ragged2e}
\usepackage{fancyhdr}
\usepackage{fontspec}
\setmainfont{Carlito}
\setsansfont{Carlito}
\newfontfamily\latinfont{Carlito}
\renewcommand{\familydefault}{\sfdefault}
\definecolor{titleblue}{HTML}{0F2F4A}
\definecolor{textblack}{HTML}{1A1A1A}
\definecolor{rule}{HTML}{C9C4C0}
\definecolor{muted}{HTML}{6B635E}
\setlength{\parindent}{0pt}\setlength{\parskip}{0.5em}

\setlength{\headheight}{28pt}\setlength{\headsep}{6pt}
\newcommand{\idmfooter}{%
  \begin{minipage}{\textwidth}
    {\color{rule}\hrule height 0.7pt}\vspace{4pt}
    \noindent
    \begin{minipage}[c]{0.66\textwidth}{\scriptsize\color{muted}Developed by the Water and Climate Lab, IIT Gandhinagar, using hydro-meteorological data from partner agencies.}\end{minipage}\hfill
    \begin{minipage}[c]{0.32\textwidth}\raggedleft
      
\includegraphics[height=0.62cm]{wcl.png}\hspace{6pt}%
      
\includegraphics[height=0.62cm]{iitgn.png}
    \end{minipage}\\[3pt]
    {\scriptsize\color{muted}© 2026 Water and Climate Lab · Indian Institute of Technology Gandhinagar · For research and demonstration purposes.}
  \end{minipage}%
}
\newcommand{\idmrunhead}{%
  \begin{minipage}{\textwidth}
    {\small\bfseries\color{titleblue}National Drought Summary}\hfill{\small\color{muted}Week ending June 3, 2026}\\[2pt]
    {\color{rule}\hrule height 0.6pt}
  \end{minipage}%
}
\fancypagestyle{idmfoot}{\fancyhf{}\renewcommand{\headrulewidth}{0pt}\renewcommand{\footrulewidth}{0pt}\fancyfoot[C]{\idmfooter}}
\fancypagestyle{idmrun}{\fancyhf{}\renewcommand{\headrulewidth}{0pt}\renewcommand{\footrulewidth}{0pt}\fancyhead[C]{\idmrunhead}\fancyfoot[C]{\idmfooter}}
\pagestyle{idmrun}

\begin{document}

\thispagestyle{idmfoot}
\noindent
\begin{minipage}[c]{0.16\textwidth}
\includegraphics[width=\linewidth]{wcl.png}\end{minipage}\hfill
\begin{minipage}[c]{0.80\textwidth}\raggedleft
  {\fontsize{19}{22}\selectfont\bfseries\color{titleblue}National Drought Summary}\\[2pt]
  {\normalsize\color{muted}Week ending June 3, 2026}
\end{minipage}
\vspace{6pt}{\color{rule}\hrule height 1pt}\vspace{10pt}
\begin{center}
  \includegraphics[width=0.60\textwidth]{cdi_map.png}\\[3pt]
  {\small\color{muted}Combined Drought Index (CDI) for the week ending June 3, 2026.}
\end{center}
\vspace{2pt}
{\large\bfseries\color{titleblue}Summary}\\[3pt]
{\justifying\normalsize\color{textblack}The national drought situation this week shows a distribution across the drought classes, with 80.4\% of the area in the Normal category and 19.6\% in the D0+ category. The total area affected by drought, specifically D0+, was 19.6\% this week, which represents a contraction of 0.5 percentage points compared to the previous week's 20.1\%. Looking further back, the area in D0+ was 15.0\% approximately one month ago. \par\medskip Gujarat exhibited the highest level of drought impact, with 85.2\% of its area in drought. Other states with significant drought presence include the NCT of Delhi, which had 69.6\% in drought, Arunanchal Pradesh at 52.8\%, Jammu \& Kashmir at 50.4\%, Mizoram at 49.8\%, and Meghalaya at 43.5\%. Conversely, Sikkim reported 0.0\% in drought, and Nagaland reported 0.7\% in drought.\par}
\vspace{8pt}
{\large\bfseries\color{titleblue}Regional Outlook}\\[2pt]
{\normalsize\color{titleblue}\textbf{North}}\enspace {\normalsize\color{textblack}Moderate, scattered drought. About 30\% of the area is in drought (D0–D4). Pockets of severe-or-worse conditions are present.} \par\smallskip {\normalsize\color{titleblue}\textbf{Northwest}}\enspace {\normalsize\color{textblack}Mostly normal, with localised dryness. About 5\% of the area is in drought (D0–D4).} \par\smallskip {\normalsize\color{titleblue}\textbf{Northeast}}\enspace {\normalsize\color{textblack}Moderate, scattered drought. About 37\% of the area is in drought (D0–D4). A notable share is in severe-or-worse categories (D2–D4).} \par\smallskip {\normalsize\color{titleblue}\textbf{East}}\enspace {\normalsize\color{textblack}Mostly normal, with localised dryness. About 11\% of the area is in drought (D0–D4). Pockets of severe-or-worse conditions are present.} \par\smallskip {\normalsize\color{titleblue}\textbf{Central}}\enspace {\normalsize\color{textblack}Mostly normal, with localised dryness. About 11\% of the area is in drought (D0–D4). Pockets of severe-or-worse conditions are present.} \par\smallskip {\normalsize\color{titleblue}\textbf{West}}\enspace {\normalsize\color{textblack}Widespread drought. About 42\% of the area is in drought (D0–D4). Pockets of severe-or-worse conditions are present.} \par\smallskip {\normalsize\color{titleblue}\textbf{South}}\enspace {\normalsize\color{textblack}Mostly normal, with localised dryness. About 8\% of the area is in drought (D0–D4).}
\end{document}
